\documentclass[a4paper,12pt,oneside]{report}
\usepackage[utf8]{inputenc}
\usepackage[english]{babel}
\usepackage{geometry}
\usepackage{amsmath}
\geometry{a4paper, margin=1in}

%% ============================================================
%% FONTS (Professional Palatino for Academic Thesis)
%% ============================================================
\usepackage{newpxtext}
\usepackage{newpxmath}
\usepackage[T1]{fontenc}

\usepackage{tocloft}
\usepackage{graphicx}
\graphicspath{{images/}} % Folder for screenshots
\usepackage{setspace}
\usepackage{xcolor}
\usepackage{float}
\usepackage{hyperref}
\usepackage{caption}
\usepackage{longtable}
\usepackage{booktabs}
\usepackage{array}
\usepackage{listings}
\usepackage{cite}
\usepackage{ragged2e}
\usepackage{pdfpages}
\usepackage{indentfirst}
\usepackage{etoolbox}
\usepackage{fancyhdr}

%% ============================================================
%% CONFIGURATION & COLORS
%% ============================================================
\definecolor{bitblue}{RGB}{0, 70, 127}
\definecolor{codegray}{rgb}{0.95,0.95,0.95}

\lstset{
  backgroundcolor=\color{codegray},
  basicstyle=\ttfamily\small,
  breaklines=true,
  frame=single,
  language=Python,
  morekeywords={await, async},
  keywordstyle=\color{bitblue}\bfseries,
  commentstyle=\color{gray}\itshape
}

\hypersetup{
    colorlinks=true,
    linkcolor=black,
    citecolor=blue,
    urlcolor=cyan
}

\onehalfspacing

%% ============================================================
%% HEADERS & FOOTERS (Promotion 2023-2026)
%% ============================================================
\setlength{\headheight}{13.6pt}
\fancyhf{}
\pagestyle{fancy}
\fancyhead[L]{\small\textit{Digital System for Prehospital Triage (ESI)}}
\fancyhead[R]{\small\textit{SAMU Burkina Faso}}
\renewcommand{\headrulewidth}{0.4pt}
\fancyfoot[L]{\small\textbf{BATIONO Jonathan} \quad Promotion: 2023--2026}
\fancyfoot[R]{\small Page \thepage}
\renewcommand{\footrulewidth}{0.4pt}

\renewcommand{\thepage}{\roman{page}}
\setcounter{tocdepth}{2}

\renewcommand{\cfttoctitlefont}{\hfil\Huge\bfseries}
\renewcommand{\cftloftitlefont}{\hfil\Huge\bfseries}
\renewcommand{\cftlottitlefont}{\hfil\Huge\bfseries}
\renewcommand{\cftbeforechapskip}{10pt}
\renewcommand{\cftbeforesecskip}{5pt}
\renewcommand{\cftbeforesubsecskip}{3pt}

\begin{document}

%% ============================================================
%% FRONT MATTER
%% ============================================================
\includepdf[pages=1]{cover.pdf}

\newpage
\thispagestyle{fancy}
\section*{Dedication}
\addcontentsline{toc}{section}{Dedication}
\vspace*{3cm}
\justifying
I dedicate this work to my family, whose unwavering support and encouragement have been the cornerstone of my academic journey. Their sacrifices and belief in my potential have provided the motivation necessary to reach this significant milestone. To my friends and mentors who walked alongside me, sharing knowledge and inspiration, this achievement is a testament to our collective determination. Finally, I dedicate this effort to the healthcare professionals of the SAMU Burkina Faso, whose dedication to saving lives inspired the creation of this digital assistance system.
\clearpage

\newpage
\thispagestyle{fancy}
\section*{Acknowledgements}
\addcontentsline{toc}{section}{Acknowledgements}
\vspace{1cm}
\justifying
I would like to express my deepest gratitude to my thesis supervisor for his invaluable guidance, patience, and constructive feedback throughout the development of this project. His expertise in software engineering and commitment to academic excellence were essential in shaping the final outcome of this work.

I also extend my sincere thanks to the administration and faculty of the Burkina Institute of Technology (BIT) for providing an environment conducive to learning and innovation. The technical and professional skills acquired during my years of study have been fundamental to the realization of this project.

My heartfelt thanks go to the staff of SAMU Burkina Faso for their cooperation and the insights provided into the challenges of pre-hospital triage. Their expert perspective ensured that the developed system remains grounded in real-world clinical needs.

Finally, I am profoundly grateful to my parents and siblings for their constant love and moral support. Without their presence, this journey would have been far more difficult.
\clearpage

\newpage
\thispagestyle{fancy}
\section*{Preface}
\addcontentsline{toc}{section}{Preface}
\justifying
The Burkina Institute of Technology (BIT), inaugurated in October 2018 in Koudougou, represents a modern vision for education in West Africa. Founded by Susanne PERTL, the institute aims to cultivate a new generation of leaders equipped with both technical expertise and an entrepreneurial mindset. By conducting classes in English and fostering international collaboration, BIT prepares its students for the global challenges of the digital age.

The institute offers specialized programs in Computer Science, Electrical Engineering, and Mechanical Engineering, complemented by a Master's degree in Artificial Intelligence. At the conclusion of their three-year undergraduate program, students are required to conduct a technical project and defend a thesis. It is within this framework of academic excellence and real-world application that the present work, focusing on the digitization of the Emergency Severity Index (ESI) triage protocol for the SAMU of Burkina Faso, was developed. This project seeks to bridge the gap between clinical needs and technological solutions to improve emergency healthcare outcomes.
\clearpage

%% ============================================================
%% ABSTRACT & RÉSUMÉ
%% ============================================================
\newpage
\thispagestyle{fancy}
\section*{Abstract}
\addcontentsline{toc}{section}{Abstract}
\justifying
The management of pre-hospital medical emergencies in Burkina Faso, coordinated by the SAMU (Service d'Aide Médicale Urgente), faces significant challenges due to the manual nature of patient triage. Triage—the critical process of prioritizing patients based on clinical severity—currently relies on verbal assessments and human experience, often leading to inconsistencies or delays in high-stress environments. This thesis explores the design and implementation of a digital decision support system specifically tailored for the SAMU, utilizing the internationally validated Emergency Severity Index (ESI) algorithm.

The proposed solution implements a robust 5-level prioritization logic through a modern web application built with a FastAPI backend and a React.js frontend. By automating the ESI decision tree, the system provides standardized guidance to operators, ensuring objective sorting of patients. Beyond the triage algorithm, the application integrates advanced security measures, supervisor dashboards for resource tracking, and automatic report generation. Functional testing confirms that the system effectively minimizes ambiguity in patient assessment, optimizing the deployment of mobile medical units (SMUR) and enhancing the overall efficiency of emergency responses in the Burkinabè context.

\vspace{0.5cm}
\noindent\textbf{Keywords:} Emergency Severity Index (ESI), SAMU, Pre-hospital Triage, Decision Support, Health Information Systems, Burkina Faso.
\clearpage

\newpage
\thispagestyle{fancy}
\section*{R\'esum\'e}
\addcontentsline{toc}{section}{R\'esum\'e}
\justifying
La régulation des urgences pré-hospitalières au Burkina Faso, assurée par le SAMU (Service d'Aide Médicale Urgente), se heurte aujourd'hui à des limites opérationnelles liées au caractère manuel du triage des patients. Ce processus, déterminant pour la survie des patients, repose principalement sur l'expérience individuelle des agents, ce qui peut engendrer des erreurs d'appréciation face à la pression ou à des informations incomplètes. Ce travail de fin de cycle porte sur la conception et le développement d'un système numérique d'aide à la décision basé sur le protocole Emergency Severity Index (ESI) à cinq niveaux.

S'appuyant sur une architecture moderne combinant FastAPI pour le backend et React.js pour le frontend, la solution permet d'automatiser l'arbre de décision ESI afin de standardiser la priorisation des patients. Le système intègre également des fonctionnalités de gestion sécurisée des accès (RBAC), un tableau de bord pour le suivi statistique en temps réel ainsi que la génération automatique de rapports d'intervention. Les résultats des tests démontrent que la numérisation du triage réduit l'ambiguïté décisionnelle et optimise l'engagement des unités mobiles de réanimation (SMUR), améliorant ainsi la réactivité de la chaîne de secours au Burkina Faso.

\vspace{0.5cm}
\noindent\textbf{Mots-cl\'es :} Emergency Severity Index (ESI), SAMU, Triage Pr\'e-hospitalier, Aide à la Décision, Systèmes d'Information en Santé, Burkina Faso.
\clearpage

%% ============================================================
%% TOC / LOF / LOT
%% ============================================================
\newpage
\thispagestyle{fancy}
\begingroup
\tableofcontents
\endgroup
\clearpage

\newpage
\thispagestyle{fancy}
\addcontentsline{toc}{section}{List of Figures}
\listoffigures
\clearpage

\newpage
\thispagestyle{fancy}
\addcontentsline{toc}{section}{List of Tables}
\listoftables
\clearpage

%% ============================================================
%% LIST OF ACRONYMS
%% ============================================================
\newpage
\thispagestyle{fancy}
\section*{List of Acronyms and Abbreviations}
\addcontentsline{toc}{section}{List of Acronyms and Abbreviations}
\renewcommand{\arraystretch}{1.3}

\begin{longtable}{|p{3.5cm}|p{10.5cm}|}
\hline
\textbf{Acronym} & \textbf{Definition} \\
\hline
API      & Application Programming Interface \\
\hline
BIT      & Burkina Institute of Technology \\
\hline
CORS     & Cross-Origin Resource Sharing \\
\hline
CSP      & Content Security Policy \\
\hline
ESI      & Emergency Severity Index \\
\hline
HTTP     & HyperText Transfer Protocol \\
\hline
JSON     & JavaScript Object Notation \\
\hline
JWT      & JSON Web Token \\
\hline
ORM      & Object-Relational Mapping \\
\hline
RBAC     & Role-Based Access Control \\
\hline
REST     & Representational State Transfer \\
\hline
SAMU     & Service d'Aide M\'edicale Urgente \\
\hline
SMUR     & Service Mobile d'Urgence et de R\'eanimation \\
\hline
SQL      & Structured Query Language \\
\hline
UML      & Unified Modeling Language \\
\hline
\end{longtable}
\clearpage

%% ============================================================
%% MAIN BODY
%% ============================================================
\renewcommand{\thepage}{\arabic{page}}
\setcounter{page}{1}

\chapter*{General Introduction}
\addcontentsline{toc}{chapter}{General Introduction}
\thispagestyle{fancy}
\justifying

The effectiveness of emergency medical systems is a critical determinant of public health outcomes, particularly in developing nations like Burkina Faso. At the heart of these systems lies the SAMU (Service d'Aide M\'edicale Urgente), an institution dedicated to providing life-saving pre-hospital care. However, the success of any emergency response is profoundly dependent on the initial sorting and prioritization of patients, a process formally known as triage. Historically, this essential task has been performed manually within the Burkinab\`e context, relying on verbal communication and the subjective experience of healthcare providers. While this traditional approach has served its purpose, it faces significant limitations in an era of increasing emergency volumes and complex clinical scenarios.

\section*{Problem Statement}
\addcontentsline{toc}{section}{Problem Statement}
In the current operational framework of the SAMU Burkina Faso, triage decisions are frequently made under extreme pressure, often based on incomplete or imprecise information provided by callers in a state of distress. This reliance on human intuition without standardized digital assistance introduces several vulnerabilities. Geographical constraints, such as urban traffic and poor road conditions in rural areas, already impose inherent delays on interventions. When these are compounded by potential inconsistencies in patient prioritization, the risk to life increases significantly. The absence of a structured, digital decision support tool means that the most critical cases may not always be identified with the necessary speed and precision, leading to suboptimal allocation of the SAMU's specialized mobile medical units (SMUR). Consequently, there is an urgent need to modernize this process through a technological solution capable of formalizing and securing triage decisions.

\section*{Objectives}
\addcontentsline{toc}{section}{Objectives}
The primary goal of this research is to design and develop a comprehensive digital decision support system for pre-hospital triage at the SAMU of Burkina Faso. This system is fundamentally based on the internationally validated Emergency Severity Index (ESI) algorithm, which provides a clinically rigorous framework for patient sorting. To achieve this, the project focuses on several specific objectives. First, we aim to analyze and define the functional requirements specific to the Burkinab\`e pre-hospital environment. Second, the project involves implementing an automated decision tree that digitizes the 5-level ESI protocol. Third, we seek to develop a secure, high-performance full-stack web application that integrates these functionalities for both frontline operators and medical supervisors. Finally, our objective includes providing advanced dashboarding and reporting features to ensure real-time visibility and auditability of triage activities.

\section*{Methodology}
\addcontentsline{toc}{section}{Methodology}
To ensure the successful realization of this project, we have adopted the Two-Track Unified Process (2TUP) as our primary development methodology. This approach is particularly well-suited for complex software projects as it allows for the parallel management of functional and technical requirements. The functional track of our work focuses on the clinical logic of the ESI algorithm and the specific workflows of the SAMU agents. Simultaneously, the technical track addresses the implementation of a modern architecture, utilizing FastAPI for the backend and React.js for the frontend, while ensuring robust data security through PostgreSQL and advanced authentication mechanisms. The convergence of these two tracks results in a solution that is both clinically accurate and technically resilient.

\section*{Thesis Organization}
\addcontentsline{toc}{section}{Thesis Organization}
This thesis is structured into two main parts, each addressing a critical dimension of the project. \textbf{Part I} establishes the theoretical foundations and reviews the state of the art in the field of digital health and emergency triage. Chapter 1 explores the essential concepts of emergency medical systems and the ESI model, while Chapter 2 details the methodology and the conceptual approach adopted. \textbf{Part II} focuses on the design and implementation of the proposed solution. Chapter 3 presents the technical modelling and the development process, and Chapter 4 analyzes the results, functional performance, and future perspectives for the system. A general conclusion synthesizes our contributions and outlines further improvements.

\clearpage


\part{Theoretical Foundations and State of the Art}
\chapter{General Information and Key Concepts}
\thispagestyle{fancy}
\justifying

This chapter provides the theoretical framework for our study, exploring the critical landscape of emergency medical systems and the specific mechanisms of patient triage. To understand the context of the proposed digital solution, it is essential to first define the roles of institutional relief organizations in Burkina Faso and the internationally validated protocols used to manage the influx of patients into emergency departments and pre-hospital care.

\section{Emergency Medical Services and the SAMU}
In Burkina Faso, the Service d'Aide M\'edicale Urgente (SAMU) was established as a cornerstone of the public health infrastructure to provide specialized and free medical care for critical patients before they reach the hospital environment. The mission of the SAMU is twofold: first, to coordinate medical responses through a centralized regulation center; and second, to deploy medicalized emergency units known as SMUR (Service Mobile d'Urgence et de R\'eanimation). These units are equipped with advanced medical technology and staffed by specialized personnel capable of performing life-saving procedures in the field. Consequently, the effectiveness of the entire emergency chain depends heavily on the accuracy of the regulation center, where the first decisions about patient priority are made.

\section{The Strategic Importance of Triage}
Medical triage is more than a simple sorting procedure; it is a strategic management tool designed to optimize limited healthcare resources and ensure the highest possible survival rates during emergencies or mass-casualty incidents. Historically, the triage concept evolved from military medicine and has been adapted for modern hospital and pre-hospital environments. The fundamental principle revolves around the idea that patients should be treated not only based on their arrival time but, more importantly, on the severity of their condition and their need for immediate intervention. In a high-pressure environment such as the SAMU regulation center, where callers provide varying levels of information, a standardized triage protocol serves as a common language between dispatchers and field teams.

\section{The Emergency Severity Index (ESI)}
To address the need for a reliable and objective triage framework, the Emergency Severity Index (ESI) was developed as a five-level clinical triage tool. Unlike traditional three-level systems, the ESI provides a more granular stratification by considering both patient acuity and the expected consumption of medical resources. The ESI algorithm is structured around four main decision points, leading to a priority ranking from 1 (immediate life threat) to 5 (least urgent). For instance, patients in cardiac arrest or with severe respiratory distress are categorized as ESI Level 1, requiring immediate resuscitation. Conversely, patients who are stable and require only simple diagnostic tests are assigned to higher levels. This objectivity is essential for minimizing the risk of clinical oversight and for standardizing care across different medical teams.

\section{Digital Decision Support in Modern Healthcare}
The integration of digital technology into clinical workflows has revolutionized the way medical decisions are made. A Digital Decision Support System (DDSS) for triage acts as an intelligent intermediary that guides the operator through the ESI algorithm, ensuring that no critical diagnostic criteria are omitted. By digitizing a clinical protocol, we can significantly reduce the cognitive load on healthcare workers, mitigate the effects of fatigue or stress, and provide an auditable record of each decision. Modern systems, built on robust architectures such as FastAPI and React, offer the performance and security required to handle sensitive patient data while providing intuitive interfaces that fit naturally into the fast-paced environment of an emergency regulation center.

\clearpage

\chapter{Methodology and Conceptual Approach}
\thispagestyle{fancy}
\justifying

This chapter discusses the methodology and the conceptual framework for our project, providing a roadmap for translating clinical needs into a functional digital system. A systematic approach to requirements gathering and software development was necessary to ensure that the developed triage tool effectively addresses the challenges faced by SAMU agents in Burkina Faso.

\section{Analysis of System Requirements}
The design of the SAMU digital triage system began with an in-depth analysis of its core functionalities and technical constraints. This involved identifying both functional requirements, which define what the system must do, and non-functional requirements, which specify how the system should operate.

\subsection{Functional Requirements}
The primary objective of the system is to assist operators in the rapid and standardized categorization of patients. This requires a digital implementation of the ESI algorithm that process vital signs and clinical criteria to determine a priority level from 1 to 5. Furthermore, the system must support secure user authentication for operators and supervisors, as well as the ability to manage and audit intervention histories. An essential feature is the supervisor's ability to oversee a real-time dashboard reflecting current triage activities and resource allocations.

\subsection{Non-Functional Requirements}
To be effective in a pre-hospital regulation center, the system must satisfy several non-functional constraints. High availability and performance are paramount, as any delay in the triage process can have serious clinical consequences. The user interface must be intuitive and responsive across different devices, from desktop workstations to mobile tablets. Data security and integrity are also critical, particularly in the handling of patient-related information, necessitating robust authentication and audit logging mechanisms.

\section{Software Development Methodology: The 2TUP Model}
To manage the complexity of this project, we have adopted the Two-Track Unified Process (2TUP) as our primary development methodology. This model is characterized by its division into two parallel tracks: a functional track and a technical track. In our project, the functional track involved the detailed analysis of the ESI protocol and the specific triage workflows within the SAMU. Simultaneously, the technical track focused on the selection and configuration of our technology stack, including the backend framework, database management, and security protocols. This dual-layered approach allowed for a specialized focus on each dimension before their convergence in the final system design and implementation.

\section{Conceptual Architecture and Data Flow}
The conceptual architecture of the SAMU triage system is based on a three-tier model, which ensures a clean separation between the user interface, the business logic, and the persistent storage of data.

\subsection{Presentation, Application, and Data Layers}
The presentation layer consists of a React.js web application that provides an interactive and responsive experience for SAMU agents. The application layer is a FastAPI-based REST API that processes triage data, executes the ESI logic, and handles user authentication. Finally, the data layer uses a PostgreSQL database to securely store operator accounts, patient records, and audit logs. This layered architecture not only improves maintainability but also ensures the scalability of the system.

\subsection{Digital Triage Workflow}
The workflow begins when an operator initiates a new triage case, entering initial patient information and clinical observations. These data are sent to the backend, where the ESI algorithm is applied in a multi-step process. Once the priority level is determined, the result is displayed back to the operator with specific recommendations, and the entire intervention is persisted in the database for later review by supervisors. This digitized process ensures that every triage decision is consistent, standardized, and auditable.

\clearpage


\part{Design and Implementation}
\chapter{Modelling and Realization Technique}
\thispagestyle{fancy}
\justifying

This chapter presents the modelling and realization of the SAMU digital triage system, providing a technical overview of its design and implementation. The process involved translating clinical workflows into a clear software architecture using Universal Modelling Language (UML) and selecting a high-performance technology stack that meets the demanding requirements of emergency medical regulation.

\section{System Modelling and UML Specifications}
The modelling phase of the SAMU triage project was essential for defining the roles of users, the structure of the data, and the dynamic behavior of the system. We developed several UML diagrams as a blueprint for development.

\subsection{Actor Interaction and Use Case Diagram}
The triage system identifies two primary actors and a supervisor. The \textbf{Operator} is the main user who performs patient triage, inputs clinical data, and views ESI results. The \textbf{Supervisor} can manage operator accounts, oversee current triage activities, and generate reports. The system automates the triage process, acting as an intelligent intermediary that applies the ESI logic to the operator's input.

\subsection{Conceptual Data Model and Class Diagram}
To ensure data integrity and robust storage, we identified several essential entities. The \textbf{Operator} class manages user accounts and security, while the \textbf{TriageCase} class stores the vital signs, clinical symptoms, and the resulting ESI level for each intervention. These entities are linked to an \textbf{AuditLog} system that records all sensitive administrative and clinical actions, ensuring transparency and accountability.

\subsection{Dynamic Behavior and Sequence Diagram}
The sequence of events during a triage session follows a clear pattern. First, the operator authenticates through a secure login interface. Second, a new triage session is initiated where clinical data is sent to the backend API. Third, the FastAPI server applies the ESI algorithm, stores the result in a PostgreSQL database, and returns the ESI level to the frontend for real-time display.

\section{Realization and Technology Stack Selection}
The technical implementation of the SAMU triage system leverages several modern technologies, selected for their reliability, performance, and security features.

\subsection{Premium Frontend with React.js}
The user interface is built with React.js and TailwindCSS, offering a premium and responsive experience that fits the institutional identity of SAMU Burkina Faso. We focused on creating a clean, modern design that minimizes cognitive load for operators during high-stress triage sessions.

\begin{figure}[H]
\centering
\includegraphics[width=0.9\textwidth]{login_page.png}
\caption{Premium Login Interface for SAMU Triage Portal}
\label{fig:login}
\end{figure}

\subsection{High-Performance Backend with FastAPI}
The application layer is implemented using FastAPI, a modern Python framework known for its high performance and native support for asynchronous programming. This choice was motivated by the need for low-latency responses during patient prioritization, as well as the ease of building secure RESTful APIs.

\begin{figure}[H]
\centering
\includegraphics[width=0.95\textwidth]{dashboard_main.png}
\caption{Supervisor Dashboard with Recharts Visualizations and 24h Reports}
\label{fig:dashboard}
\end{figure}

\subsection{Robust Data Management with PostgreSQL}
To move from local SQLite tests to a production-ready environment, we migrated the system to a PostgreSQL database. This migration ensures robust transactional support, efficient handling of historical data, and a secure platform for auditing and reporting.

\begin{figure}[H]
\centering
\includegraphics[width=0.9\textwidth]{user_management.png}
\caption{Admin Interface for Operator and Specialist Management}
\label{fig:admin}
\end{figure}

\section{Security Architecture and Implementation}
Security was a primary consideration throughout the project, culminating in several hardening measures. We implemented secure authentication using JSON Web Tokens (JWT) and Bcrypt for password hashing. Furthermore, administrative actions are protected by a role-based access control (RBAC) system, ensuring that only supervisors can access sensitive user management features. Additional security headers, such as Content-Security-Policy (CSP) and brute-force protection through login delays and lockouts, were also integrated to protect the portal from external threats.

\clearpage

\chapter{Results, Analysis and Discussion}
\thispagestyle{fancy}
\justifying

This chapter presents the results of the SAMU digital triage system development and provides an analysis of its clinical and technical performance. We successfully implemented the ESI 5-level protocol through an interactive web-based decision tree and validated its effectiveness through several functional test scenarios.

\section{System Implementation and Functional Results}
The primary outcome of the project is a fully functional web portal that enables SAMU operators to prioritize patients in a standardized and efficient manner. The user experience is designed to be as simple as possible, allowing for rapid data entry without losing essential clinical precision.

\subsection{Standardized ESI Triage Process}
The core of the system is the interactive decision tree, which follows the ESI v4 algorithm logic. Users walk through a multi-step form to determine the appropriate priority level for each patient. This process ensures that every critical diagnostic question is answered, minimizing the risk of misclassification.

\begin{figure}[H]
\centering
\includegraphics[width=0.9\textwidth]{triage_multi_step.png}
\caption{Step-by-Step Triage Form for SAMU Operators}
\label{fig:triage_form}
\end{figure}

\subsection{Automated ESI Level Calculation}
Once the clinical data is processed, the system returns a clear ESI level (1 to 5) with specific recommendations for each priority. For example, an ESI Level 1 result is accompanied by immediate resuscitation warnings and priority unit dispatch alerts.

\begin{figure}[H]
\centering
\includegraphics[width=0.8\textwidth]{esi_result_box.png}
\caption{ESI Triage Result with Level and Immediate Priority Indicator}
\label{fig:esi_result}
\end{figure}

\subsection{Intervention History and PDF Reporting}
A critical requirement for the SAMU is the ability to record and archive every triage decision. The system provides a historical view of all interventions and allows operators to export a standardized PDF report for each session. This feature is essential for clinical auditability and for maintaining a permanent medical record.

\begin{figure}[H]
\centering
\includegraphics[width=0.85\textwidth]{pdf_report_sample.png}
\caption{Exported PDF Triage Report with Standardized SAMU Formatting}
\label{fig:pdf_report}
\end{figure}

\section{Performance Evaluation and Testing}
To ensure the reliability of the system, we conducted several functional tests across different ESI scenarios. We compared the system's output with manual expert triage results to verify accuracy.

\subsection{Functional Test Scenarios}
In tests involving Level 1 (Immediate Life Threat) and Level 2 (High Risk) scenarios, the system achieved a 100\% match with expected clinical priorities. For ESI Levels 3 to 5, where resource counting is required, the system consistently applied the correct logic, demonstrating its value in eliminating ambiguity.

\begin{table}[H]
\centering
\caption{Functional Validation of ESI Levels across Test Scenarios}
\begin{tabular}{|l|l|l|l|}
\hline
\textbf{Scenario} & \textbf{Clinical Condition} & \textbf{Expected ESI} & \textbf{System Result} \\
\hline
Scenario A & Cardiac Arrest / Apnea      & Level 1      & Level 1 (Success) \\
\hline
Scenario B & Acute Chest Pain / High Risk & Level 2      & Level 2 (Success) \\
\hline
Scenario C & Minor Fracture (2 resources) & Level 3      & Level 3 (Success) \\
\hline
Scenario D & Simple Laceration (1 res.) & Level 4      & Level 4 (Success) \\
\hline
Scenario E & Prescription Refill (0 res.)& Level 5      & Level 5 (Success) \\
\hline
\end{tabular}
\label{tab:functional_tests}
\end{table}

\section{Analysis and Discussion}
The development of the SAMU digital triage portal represents a significant step towards the modernization of the pre-hospital emergency chain in Burkina Faso. Our analysis focuses on both the technical achievements and the potential clinical impact of the system.

\subsection{Clinical and Operational Impact}
The system successfully addresses the limitations of manual triage by providing an objective, auditable framework for patient sorting. This standardization is particularly important in high-volume situations or when dealing with inexperienced staff. The automated 24h reports and supervisor dashboards provide SAMU management with valuable data for resource planning and quality improvement.

\subsection{Technical Considerations and System Hardening}
From a technical standpoint, the choice of FastAPI and React has proven to be effective in delivering a high-performance and secure system. The migration to PostgreSQL and the implementation of advanced security headers ensure that the application is production-ready. The system's modular design allows for future extensions, such as integration with hospital information systems or mobile dispatch units.

\subsection{Limitations and Future Perspectives}
While the system is fully functional, certain limitations must be acknowledged. The current version requires a reliable network connection to communicate with the central API. Future versions could incorporate offline capabilities or mobile-only interfaces for use in remote rural areas. Furthermore, the integration of real-time GPS tracking for SMUR units would further enhance the system's ability to coordinate large-scale emergency responses.

\clearpage


\chapter*{General Conclusion and Perspectives}
\addcontentsline{toc}{chapter}{General Conclusion and Perspectives}
\thispagestyle{fancy}
\justifying

The design and development of a digital decision support system for pre-hospital triage at the SAMU of Burkina Faso, based on the Emergency Severity Index (ESI) protocol, represents a fundamental achievement in the modernization of our emergency medical infrastructure. Throughout this project, we have transformed a manual and subjective process into a standardized, efficient, and secure digital platform that addresses the real-world challenges faced by SAMU operators and supervisors.

\section*{Synthesis of Contributions}
By adopting the Two-Track Unified Process (2TUP) methodology, we successfully navigated both the complex clinical logic of the 5-level ESI algorithm and the technical demands of a modern full-stack web application. The developed system provides a robust solution for patient prioritization, integrating real-time ESI level calculation with comprehensive audit logging and reporting features. Our functional testing and performance analysis have demonstrated that this digital approach significantly reduces ambiguity in patient categorization, ultimately optimizing the deployment of mobile medical units (SMUR). Furthermore, the integration of advanced security measures and intuitive dashboards ensures that the SAMU triage portal is production-ready and capable of handling sensitive medical information with the highest standards of integrity.

\section*{Key Perspectives for Further Development}
While the current version of the system provides a solid foundation for digitized triage, several avenues for future research and development have been identified. From a technical standpoint, the implementation of a mobile-first application with offline capabilities would further enhance the system's utility in remote or low-connectivity areas of Burkina Faso. Operationally, the integration of the triage portal with hospital emergency departments could streamline the handoff process, ensuring that receiving facilities are informed of incoming patients' priority levels in real time. Additionally, the potential for incorporating geographic information systems (GIS) for fleet tracking and optimization offers a promising direction for improving overall emergency response times.

Ultimately, the digital SAMU triage system is more than just a software tool; it is a clinical and operational asset that contributes to the core mission of saving lives by ensuring that the right care reaches the right patient at the right time.

\clearpage


%% ============================================================
%% BIBLIOGRAPHY (At the end)
%% ============================================================
\clearpage
\newpage
\thispagestyle{fancy}
\addcontentsline{toc}{chapter}{Bibliography}
\bibliographystyle{plain}
\bibliography{references}

\end{document}
